% =====================================================================
%  theme.tex --- the one visual language of the project
%
%  The palette, the axis, the legend strip, the interval bar, the
%  hatched bar, the reference band, the panel label and the table rules.
%  Every figure in the paper and every figure and table in the report
%  reads from this file, so a colour or a line weight is changed in ONE
%  place.
%
%  It needs only xcolor, tikz and pgfplots, which the paper and the
%  report both already load, and it loads the TikZ `patterns' library
%  itself, so no preamble line is strictly required.  The recommended
%  line for the paper's preamble (after \pgfplotsset{compat=1.18}) is
%
%      % =====================================================================
%  theme.tex --- the one visual language of the project
%
%  The palette, the axis, the legend strip, the interval bar, the
%  hatched bar, the reference band, the panel label and the table rules.
%  Every figure in the paper and every figure and table in the report
%  reads from this file, so a colour or a line weight is changed in ONE
%  place.
%
%  It needs only xcolor, tikz and pgfplots, which the paper and the
%  report both already load, and it loads the TikZ `patterns' library
%  itself, so no preamble line is strictly required.  The recommended
%  line for the paper's preamble (after \pgfplotsset{compat=1.18}) is
%
%      % =====================================================================
%  theme.tex --- the one visual language of the project
%
%  The palette, the axis, the legend strip, the interval bar, the
%  hatched bar, the reference band, the panel label and the table rules.
%  Every figure in the paper and every figure and table in the report
%  reads from this file, so a colour or a line weight is changed in ONE
%  place.
%
%  It needs only xcolor, tikz and pgfplots, which the paper and the
%  report both already load, and it loads the TikZ `patterns' library
%  itself, so no preamble line is strictly required.  The recommended
%  line for the paper's preamble (after \pgfplotsset{compat=1.18}) is
%
%      % =====================================================================
%  theme.tex --- the one visual language of the project
%
%  The palette, the axis, the legend strip, the interval bar, the
%  hatched bar, the reference band, the panel label and the table rules.
%  Every figure in the paper and every figure and table in the report
%  reads from this file, so a colour or a line weight is changed in ONE
%  place.
%
%  It needs only xcolor, tikz and pgfplots, which the paper and the
%  report both already load, and it loads the TikZ `patterns' library
%  itself, so no preamble line is strictly required.  The recommended
%  line for the paper's preamble (after \pgfplotsset{compat=1.18}) is
%
%      \input{ml_impl/analysis/theme.tex}
%
%  The generated figure fragments bootstrap the same file if the
%  preamble has not, so a paper that has not been edited still compiles.
% =====================================================================

% the marker that says the theme is loaded: the fragments test it, so
% the file is never read twice (a second \pgfplotsset would be harmless,
% a second \newcommand would not)
\def\uavtheme{2026-10-21}

% the diagonal hatch of the primary bar
\usetikzlibrary{patterns}


% ---------------------------------------------------------------------
% 1. the palette
%    uavblue is the project's primary; uavref is the reference grey the
%    primary is always read against; uavink is every piece of type; the
%    rest are the accents the paper's own preamble already defines.
% ---------------------------------------------------------------------
\providecolor{uavblue}{RGB}{37,99,235}
\providecolor{uavcyan}{RGB}{8,145,178}
\providecolor{uavgreen}{RGB}{22,163,74}
\providecolor{uavorange}{RGB}{234,88,12}
\providecolor{uavred}{RGB}{220,38,38}
\providecolor{uavpurple}{RGB}{124,58,237}

\providecolor{uavink}{RGB}{30,41,59}        % type and rules
\providecolor{uavmuted}{RGB}{100,116,139}   % secondary type
\providecolor{uavgrid}{RGB}{214,219,227}    % gridlines, legend frame
\providecolor{uavref}{RGB}{179,179,179}     % the reference series
\providecolor{uavband}{RGB}{233,236,241}    % a reference region
\providecolor{uavzebra}{RGB}{247,248,251}   % a shaded table row
\providecolor{uavlight}{RGB}{241,245,249}   % a table head
\providecolor{uavwhite}{RGB}{255,255,255}

% ---------------------------------------------------------------------
% 1b. THE BLUE RAMP
%     The project's figures are read against ONE hue, so a reader learns
%     the grammar once: blue is the method, grey is the reference, and
%     the ramp says "the same quantity, more or less of it".  The ramp
%     is the primary hue at five strengths, from the tint a table row
%     is filled with to the shade that carries the darkest series.
%     The accents (red, amber, green) are reserved for MEANINGS -- a
%     bound that is not met, a criterion that disagrees, an effect that
%     is resolved -- and are never used decoratively.
% ---------------------------------------------------------------------
\providecolor{uavp100}{RGB}{219,234,254}   % the tint: a highlighted row
\providecolor{uavp75}{RGB}{147,183,253}    % a pale band
\providecolor{uavp50}{RGB}{96,145,246}     % a light series
\providecolor{uavp25}{RGB}{66,115,241}     % a mid series
\providecolor{uavdeep}{RGB}{30,64,175}     % the darkest: a threshold
\providecolor{uavinkblue}{RGB}{23,44,110}   % a label on the primary

% the greys, in the same order of strength
\providecolor{uavg90}{RGB}{241,243,246}    % a panel's ground
\providecolor{uavg70}{RGB}{214,219,227}    % a gridline
\providecolor{uavg50}{RGB}{179,186,196}    % a reference series
\providecolor{uavg30}{RGB}{116,124,136}    % secondary type
\providecolor{uavg10}{RGB}{45,55,72}       % primary type

% the accents, with the meaning each one carries
\providecolor{uavaccent}{RGB}{220,38,38}    % a bound that is not met
\providecolor{uavwarn}{RGB}{234,120,12}     % a criterion that disagrees
\providecolor{uavok}{RGB}{22,163,74}       % an effect that is resolved


% ---------------------------------------------------------------------
% 2. the axis
%    one spine on the left and one at the bottom, no frame, light dashed
%    gridlines, small dark tick labels and a bold axis label.
% ---------------------------------------------------------------------
\pgfplotsset{
  uavaxis/.style={
      % width/height are the PLOT AREA, so the picture width is a linear
      % function of the axis width and the fit to the column is exact
      scale only axis,
      axis x line=bottom,
      axis y line=left,
      axis line style={uavink, line width=0.6pt},
      tick style={uavink, line width=0.6pt},
      major tick length=2.2pt,
      tick align=outside,
      tick label style={font=\scriptsize, text=uavink, inner sep=1.6pt},
      label style={font=\small\bfseries, text=uavink},
      grid style={uavgrid, dashed, line width=0.5pt},
      legend cell align=left,
      clip=false,
  },
  % the boxed legend strip that runs across the top of a figure
  uavlegend/.style={
      legend style={
          draw=uavgrid,
          fill=uavwhite,
          line width=0.5pt,
          rounded corners=1pt,
          font=\scriptsize,
          text=uavink,
          inner xsep=4pt,
          inner ysep=3pt,
          nodes={inner sep=1.5pt, text=uavink},
          cells={anchor=west},
          column sep=5pt,
          row sep=1pt,
          at={(0.5,1.045)},
          anchor=south,
          /tikz/every even column/.append style={column sep=7pt},
      },
      legend image post style={xscale=0.85},
  },
  % a light grey band: a reference region of the x axis
  uavband/.style={
      draw=none,
      fill=uavband,
      forget plot,
  },
  % a level line: the reference value the figure is read against
  uavlevel/.style={
      draw=uavink,
      line width=0.9pt,
  },
  % the interval bar of a forest plot: thick, round caps, no end marks,
  % a filled marker on the point estimate.  The direction keys cannot be
  % written with "error bars/.cd" inside a style (the path would be
  % changed for every later key of the plot), so they take the full path.
  uavinterval/.style={
      only marks,
      mark=*,
      mark size=2.4pt,
      mark options={line width=0.5pt},
      error bars/error bar style={line width=2.6pt, line cap=round},
      error bars/error mark=none,
  },
  uavintervalx/.style={
      uavinterval,
      error bars/x dir=both,
      error bars/x explicit,
  },
  uavintervaly/.style={
      uavinterval,
      error bars/y dir=both,
      error bars/y explicit,
  },
  % the primary series as a bar: filled, hatched in white, blue edge
  uavhatchbar/.style={
      ybar,
      bar width=0.62,
      fill=uavblue,
      draw=uavblue,
      line width=0.5pt,
      pattern=north east lines,
      pattern color=uavwhite,
  },
  % the reference series as a bar: solid grey, no hatch
  uavrefbar/.style={
      ybar,
      bar width=0.62,
      fill=uavref,
      draw=uavref!70!black,
      line width=0.5pt,
  },
}

\tikzset{
  % the separator between two groups of a categorical axis.  It is used
  % on \draw, which reads keys from the /tikz path, so it is registered
  % with \tikzset rather than in the pgfplots-only namespace above;
  % otherwise pgfkeys reports "/tikz/uavsep unknown" and the paper's
  % gate fails.
  uavsep/.style={
      draw=uavgrid,
      line width=0.5pt,
      dash pattern=on 2pt off 2pt,
  },
}


% ---------------------------------------------------------------------
% 3. the panel label: "(a) By skill", bold, above its panel
% ---------------------------------------------------------------------
\newcommand{\uavpanel}[2]{{\normalfont\bfseries\small (#1)~#2}}


% ---------------------------------------------------------------------
% 4. the table
%    booktabs rules in the theme's colours, and a tinted head row.
%    The colouring is skipped when colortbl is not loaded, so the same
%    macros are safe in a document that only loads plain xcolor.
% ---------------------------------------------------------------------
\newcommand{\uavtoprule}{%
    \ifdefined\arrayrulecolor\arrayrulecolor{uavink}\fi
    \specialrule{1pt}{0pt}{3pt}}
\newcommand{\uavmidrule}{%
    \ifdefined\arrayrulecolor\arrayrulecolor{uavgrid}\fi
    \specialrule{0.8pt}{2pt}{3pt}}
\newcommand{\uavbottomrule}{%
    \ifdefined\arrayrulecolor\arrayrulecolor{uavink}\fi
    \specialrule{1pt}{4pt}{0pt}}
\newcommand{\uavheadrow}{\ifdefined\rowcolor\rowcolor{uavlight}\fi}
\newcommand{\uavzebrarow}{\ifdefined\rowcolor\rowcolor{uavzebra}\fi}

%
%  The generated figure fragments bootstrap the same file if the
%  preamble has not, so a paper that has not been edited still compiles.
% =====================================================================

% the marker that says the theme is loaded: the fragments test it, so
% the file is never read twice (a second \pgfplotsset would be harmless,
% a second \newcommand would not)
\def\uavtheme{2026-10-21}

% the diagonal hatch of the primary bar
\usetikzlibrary{patterns}


% ---------------------------------------------------------------------
% 1. the palette
%    uavblue is the project's primary; uavref is the reference grey the
%    primary is always read against; uavink is every piece of type; the
%    rest are the accents the paper's own preamble already defines.
% ---------------------------------------------------------------------
\providecolor{uavblue}{RGB}{37,99,235}
\providecolor{uavcyan}{RGB}{8,145,178}
\providecolor{uavgreen}{RGB}{22,163,74}
\providecolor{uavorange}{RGB}{234,88,12}
\providecolor{uavred}{RGB}{220,38,38}
\providecolor{uavpurple}{RGB}{124,58,237}

\providecolor{uavink}{RGB}{30,41,59}        % type and rules
\providecolor{uavmuted}{RGB}{100,116,139}   % secondary type
\providecolor{uavgrid}{RGB}{214,219,227}    % gridlines, legend frame
\providecolor{uavref}{RGB}{179,179,179}     % the reference series
\providecolor{uavband}{RGB}{233,236,241}    % a reference region
\providecolor{uavzebra}{RGB}{247,248,251}   % a shaded table row
\providecolor{uavlight}{RGB}{241,245,249}   % a table head
\providecolor{uavwhite}{RGB}{255,255,255}

% ---------------------------------------------------------------------
% 1b. THE BLUE RAMP
%     The project's figures are read against ONE hue, so a reader learns
%     the grammar once: blue is the method, grey is the reference, and
%     the ramp says "the same quantity, more or less of it".  The ramp
%     is the primary hue at five strengths, from the tint a table row
%     is filled with to the shade that carries the darkest series.
%     The accents (red, amber, green) are reserved for MEANINGS -- a
%     bound that is not met, a criterion that disagrees, an effect that
%     is resolved -- and are never used decoratively.
% ---------------------------------------------------------------------
\providecolor{uavp100}{RGB}{219,234,254}   % the tint: a highlighted row
\providecolor{uavp75}{RGB}{147,183,253}    % a pale band
\providecolor{uavp50}{RGB}{96,145,246}     % a light series
\providecolor{uavp25}{RGB}{66,115,241}     % a mid series
\providecolor{uavdeep}{RGB}{30,64,175}     % the darkest: a threshold
\providecolor{uavinkblue}{RGB}{23,44,110}   % a label on the primary

% the greys, in the same order of strength
\providecolor{uavg90}{RGB}{241,243,246}    % a panel's ground
\providecolor{uavg70}{RGB}{214,219,227}    % a gridline
\providecolor{uavg50}{RGB}{179,186,196}    % a reference series
\providecolor{uavg30}{RGB}{116,124,136}    % secondary type
\providecolor{uavg10}{RGB}{45,55,72}       % primary type

% the accents, with the meaning each one carries
\providecolor{uavaccent}{RGB}{220,38,38}    % a bound that is not met
\providecolor{uavwarn}{RGB}{234,120,12}     % a criterion that disagrees
\providecolor{uavok}{RGB}{22,163,74}       % an effect that is resolved


% ---------------------------------------------------------------------
% 2. the axis
%    one spine on the left and one at the bottom, no frame, light dashed
%    gridlines, small dark tick labels and a bold axis label.
% ---------------------------------------------------------------------
\pgfplotsset{
  uavaxis/.style={
      % width/height are the PLOT AREA, so the picture width is a linear
      % function of the axis width and the fit to the column is exact
      scale only axis,
      axis x line=bottom,
      axis y line=left,
      axis line style={uavink, line width=0.6pt},
      tick style={uavink, line width=0.6pt},
      major tick length=2.2pt,
      tick align=outside,
      tick label style={font=\scriptsize, text=uavink, inner sep=1.6pt},
      label style={font=\small\bfseries, text=uavink},
      grid style={uavgrid, dashed, line width=0.5pt},
      legend cell align=left,
      clip=false,
  },
  % the boxed legend strip that runs across the top of a figure
  uavlegend/.style={
      legend style={
          draw=uavgrid,
          fill=uavwhite,
          line width=0.5pt,
          rounded corners=1pt,
          font=\scriptsize,
          text=uavink,
          inner xsep=4pt,
          inner ysep=3pt,
          nodes={inner sep=1.5pt, text=uavink},
          cells={anchor=west},
          column sep=5pt,
          row sep=1pt,
          at={(0.5,1.045)},
          anchor=south,
          /tikz/every even column/.append style={column sep=7pt},
      },
      legend image post style={xscale=0.85},
  },
  % a light grey band: a reference region of the x axis
  uavband/.style={
      draw=none,
      fill=uavband,
      forget plot,
  },
  % a level line: the reference value the figure is read against
  uavlevel/.style={
      draw=uavink,
      line width=0.9pt,
  },
  % the interval bar of a forest plot: thick, round caps, no end marks,
  % a filled marker on the point estimate.  The direction keys cannot be
  % written with "error bars/.cd" inside a style (the path would be
  % changed for every later key of the plot), so they take the full path.
  uavinterval/.style={
      only marks,
      mark=*,
      mark size=2.4pt,
      mark options={line width=0.5pt},
      error bars/error bar style={line width=2.6pt, line cap=round},
      error bars/error mark=none,
  },
  uavintervalx/.style={
      uavinterval,
      error bars/x dir=both,
      error bars/x explicit,
  },
  uavintervaly/.style={
      uavinterval,
      error bars/y dir=both,
      error bars/y explicit,
  },
  % the primary series as a bar: filled, hatched in white, blue edge
  uavhatchbar/.style={
      ybar,
      bar width=0.62,
      fill=uavblue,
      draw=uavblue,
      line width=0.5pt,
      pattern=north east lines,
      pattern color=uavwhite,
  },
  % the reference series as a bar: solid grey, no hatch
  uavrefbar/.style={
      ybar,
      bar width=0.62,
      fill=uavref,
      draw=uavref!70!black,
      line width=0.5pt,
  },
}

\tikzset{
  % the separator between two groups of a categorical axis.  It is used
  % on \draw, which reads keys from the /tikz path, so it is registered
  % with \tikzset rather than in the pgfplots-only namespace above;
  % otherwise pgfkeys reports "/tikz/uavsep unknown" and the paper's
  % gate fails.
  uavsep/.style={
      draw=uavgrid,
      line width=0.5pt,
      dash pattern=on 2pt off 2pt,
  },
}


% ---------------------------------------------------------------------
% 3. the panel label: "(a) By skill", bold, above its panel
% ---------------------------------------------------------------------
\newcommand{\uavpanel}[2]{{\normalfont\bfseries\small (#1)~#2}}


% ---------------------------------------------------------------------
% 4. the table
%    booktabs rules in the theme's colours, and a tinted head row.
%    The colouring is skipped when colortbl is not loaded, so the same
%    macros are safe in a document that only loads plain xcolor.
% ---------------------------------------------------------------------
\newcommand{\uavtoprule}{%
    \ifdefined\arrayrulecolor\arrayrulecolor{uavink}\fi
    \specialrule{1pt}{0pt}{3pt}}
\newcommand{\uavmidrule}{%
    \ifdefined\arrayrulecolor\arrayrulecolor{uavgrid}\fi
    \specialrule{0.8pt}{2pt}{3pt}}
\newcommand{\uavbottomrule}{%
    \ifdefined\arrayrulecolor\arrayrulecolor{uavink}\fi
    \specialrule{1pt}{4pt}{0pt}}
\newcommand{\uavheadrow}{\ifdefined\rowcolor\rowcolor{uavlight}\fi}
\newcommand{\uavzebrarow}{\ifdefined\rowcolor\rowcolor{uavzebra}\fi}

%
%  The generated figure fragments bootstrap the same file if the
%  preamble has not, so a paper that has not been edited still compiles.
% =====================================================================

% the marker that says the theme is loaded: the fragments test it, so
% the file is never read twice (a second \pgfplotsset would be harmless,
% a second \newcommand would not)
\def\uavtheme{2026-10-21}

% the diagonal hatch of the primary bar
\usetikzlibrary{patterns}


% ---------------------------------------------------------------------
% 1. the palette
%    uavblue is the project's primary; uavref is the reference grey the
%    primary is always read against; uavink is every piece of type; the
%    rest are the accents the paper's own preamble already defines.
% ---------------------------------------------------------------------
\providecolor{uavblue}{RGB}{37,99,235}
\providecolor{uavcyan}{RGB}{8,145,178}
\providecolor{uavgreen}{RGB}{22,163,74}
\providecolor{uavorange}{RGB}{234,88,12}
\providecolor{uavred}{RGB}{220,38,38}
\providecolor{uavpurple}{RGB}{124,58,237}

\providecolor{uavink}{RGB}{30,41,59}        % type and rules
\providecolor{uavmuted}{RGB}{100,116,139}   % secondary type
\providecolor{uavgrid}{RGB}{214,219,227}    % gridlines, legend frame
\providecolor{uavref}{RGB}{179,179,179}     % the reference series
\providecolor{uavband}{RGB}{233,236,241}    % a reference region
\providecolor{uavzebra}{RGB}{247,248,251}   % a shaded table row
\providecolor{uavlight}{RGB}{241,245,249}   % a table head
\providecolor{uavwhite}{RGB}{255,255,255}

% ---------------------------------------------------------------------
% 1b. THE BLUE RAMP
%     The project's figures are read against ONE hue, so a reader learns
%     the grammar once: blue is the method, grey is the reference, and
%     the ramp says "the same quantity, more or less of it".  The ramp
%     is the primary hue at five strengths, from the tint a table row
%     is filled with to the shade that carries the darkest series.
%     The accents (red, amber, green) are reserved for MEANINGS -- a
%     bound that is not met, a criterion that disagrees, an effect that
%     is resolved -- and are never used decoratively.
% ---------------------------------------------------------------------
\providecolor{uavp100}{RGB}{219,234,254}   % the tint: a highlighted row
\providecolor{uavp75}{RGB}{147,183,253}    % a pale band
\providecolor{uavp50}{RGB}{96,145,246}     % a light series
\providecolor{uavp25}{RGB}{66,115,241}     % a mid series
\providecolor{uavdeep}{RGB}{30,64,175}     % the darkest: a threshold
\providecolor{uavinkblue}{RGB}{23,44,110}   % a label on the primary

% the greys, in the same order of strength
\providecolor{uavg90}{RGB}{241,243,246}    % a panel's ground
\providecolor{uavg70}{RGB}{214,219,227}    % a gridline
\providecolor{uavg50}{RGB}{179,186,196}    % a reference series
\providecolor{uavg30}{RGB}{116,124,136}    % secondary type
\providecolor{uavg10}{RGB}{45,55,72}       % primary type

% the accents, with the meaning each one carries
\providecolor{uavaccent}{RGB}{220,38,38}    % a bound that is not met
\providecolor{uavwarn}{RGB}{234,120,12}     % a criterion that disagrees
\providecolor{uavok}{RGB}{22,163,74}       % an effect that is resolved


% ---------------------------------------------------------------------
% 2. the axis
%    one spine on the left and one at the bottom, no frame, light dashed
%    gridlines, small dark tick labels and a bold axis label.
% ---------------------------------------------------------------------
\pgfplotsset{
  uavaxis/.style={
      % width/height are the PLOT AREA, so the picture width is a linear
      % function of the axis width and the fit to the column is exact
      scale only axis,
      axis x line=bottom,
      axis y line=left,
      axis line style={uavink, line width=0.6pt},
      tick style={uavink, line width=0.6pt},
      major tick length=2.2pt,
      tick align=outside,
      tick label style={font=\scriptsize, text=uavink, inner sep=1.6pt},
      label style={font=\small\bfseries, text=uavink},
      grid style={uavgrid, dashed, line width=0.5pt},
      legend cell align=left,
      clip=false,
  },
  % the boxed legend strip that runs across the top of a figure
  uavlegend/.style={
      legend style={
          draw=uavgrid,
          fill=uavwhite,
          line width=0.5pt,
          rounded corners=1pt,
          font=\scriptsize,
          text=uavink,
          inner xsep=4pt,
          inner ysep=3pt,
          nodes={inner sep=1.5pt, text=uavink},
          cells={anchor=west},
          column sep=5pt,
          row sep=1pt,
          at={(0.5,1.045)},
          anchor=south,
          /tikz/every even column/.append style={column sep=7pt},
      },
      legend image post style={xscale=0.85},
  },
  % a light grey band: a reference region of the x axis
  uavband/.style={
      draw=none,
      fill=uavband,
      forget plot,
  },
  % a level line: the reference value the figure is read against
  uavlevel/.style={
      draw=uavink,
      line width=0.9pt,
  },
  % the interval bar of a forest plot: thick, round caps, no end marks,
  % a filled marker on the point estimate.  The direction keys cannot be
  % written with "error bars/.cd" inside a style (the path would be
  % changed for every later key of the plot), so they take the full path.
  uavinterval/.style={
      only marks,
      mark=*,
      mark size=2.4pt,
      mark options={line width=0.5pt},
      error bars/error bar style={line width=2.6pt, line cap=round},
      error bars/error mark=none,
  },
  uavintervalx/.style={
      uavinterval,
      error bars/x dir=both,
      error bars/x explicit,
  },
  uavintervaly/.style={
      uavinterval,
      error bars/y dir=both,
      error bars/y explicit,
  },
  % the primary series as a bar: filled, hatched in white, blue edge
  uavhatchbar/.style={
      ybar,
      bar width=0.62,
      fill=uavblue,
      draw=uavblue,
      line width=0.5pt,
      pattern=north east lines,
      pattern color=uavwhite,
  },
  % the reference series as a bar: solid grey, no hatch
  uavrefbar/.style={
      ybar,
      bar width=0.62,
      fill=uavref,
      draw=uavref!70!black,
      line width=0.5pt,
  },
}

\tikzset{
  % the separator between two groups of a categorical axis.  It is used
  % on \draw, which reads keys from the /tikz path, so it is registered
  % with \tikzset rather than in the pgfplots-only namespace above;
  % otherwise pgfkeys reports "/tikz/uavsep unknown" and the paper's
  % gate fails.
  uavsep/.style={
      draw=uavgrid,
      line width=0.5pt,
      dash pattern=on 2pt off 2pt,
  },
}


% ---------------------------------------------------------------------
% 3. the panel label: "(a) By skill", bold, above its panel
% ---------------------------------------------------------------------
\newcommand{\uavpanel}[2]{{\normalfont\bfseries\small (#1)~#2}}


% ---------------------------------------------------------------------
% 4. the table
%    booktabs rules in the theme's colours, and a tinted head row.
%    The colouring is skipped when colortbl is not loaded, so the same
%    macros are safe in a document that only loads plain xcolor.
% ---------------------------------------------------------------------
\newcommand{\uavtoprule}{%
    \ifdefined\arrayrulecolor\arrayrulecolor{uavink}\fi
    \specialrule{1pt}{0pt}{3pt}}
\newcommand{\uavmidrule}{%
    \ifdefined\arrayrulecolor\arrayrulecolor{uavgrid}\fi
    \specialrule{0.8pt}{2pt}{3pt}}
\newcommand{\uavbottomrule}{%
    \ifdefined\arrayrulecolor\arrayrulecolor{uavink}\fi
    \specialrule{1pt}{4pt}{0pt}}
\newcommand{\uavheadrow}{\ifdefined\rowcolor\rowcolor{uavlight}\fi}
\newcommand{\uavzebrarow}{\ifdefined\rowcolor\rowcolor{uavzebra}\fi}

%
%  The generated figure fragments bootstrap the same file if the
%  preamble has not, so a paper that has not been edited still compiles.
% =====================================================================

% the marker that says the theme is loaded: the fragments test it, so
% the file is never read twice (a second \pgfplotsset would be harmless,
% a second \newcommand would not)
\def\uavtheme{2026-10-21}

% the diagonal hatch of the primary bar
\usetikzlibrary{patterns}


% ---------------------------------------------------------------------
% 1. the palette
%    uavblue is the project's primary; uavref is the reference grey the
%    primary is always read against; uavink is every piece of type; the
%    rest are the accents the paper's own preamble already defines.
% ---------------------------------------------------------------------
\providecolor{uavblue}{RGB}{37,99,235}
\providecolor{uavcyan}{RGB}{8,145,178}
\providecolor{uavgreen}{RGB}{22,163,74}
\providecolor{uavorange}{RGB}{234,88,12}
\providecolor{uavred}{RGB}{220,38,38}
\providecolor{uavpurple}{RGB}{124,58,237}

\providecolor{uavink}{RGB}{30,41,59}        % type and rules
\providecolor{uavmuted}{RGB}{100,116,139}   % secondary type
\providecolor{uavgrid}{RGB}{214,219,227}    % gridlines, legend frame
\providecolor{uavref}{RGB}{179,179,179}     % the reference series
\providecolor{uavband}{RGB}{233,236,241}    % a reference region
\providecolor{uavzebra}{RGB}{247,248,251}   % a shaded table row
\providecolor{uavlight}{RGB}{241,245,249}   % a table head
\providecolor{uavwhite}{RGB}{255,255,255}

% ---------------------------------------------------------------------
% 1b. THE BLUE RAMP
%     The project's figures are read against ONE hue, so a reader learns
%     the grammar once: blue is the method, grey is the reference, and
%     the ramp says "the same quantity, more or less of it".  The ramp
%     is the primary hue at five strengths, from the tint a table row
%     is filled with to the shade that carries the darkest series.
%     The accents (red, amber, green) are reserved for MEANINGS -- a
%     bound that is not met, a criterion that disagrees, an effect that
%     is resolved -- and are never used decoratively.
% ---------------------------------------------------------------------
\providecolor{uavp100}{RGB}{219,234,254}   % the tint: a highlighted row
\providecolor{uavp75}{RGB}{147,183,253}    % a pale band
\providecolor{uavp50}{RGB}{96,145,246}     % a light series
\providecolor{uavp25}{RGB}{66,115,241}     % a mid series
\providecolor{uavdeep}{RGB}{30,64,175}     % the darkest: a threshold
\providecolor{uavinkblue}{RGB}{23,44,110}   % a label on the primary

% the greys, in the same order of strength
\providecolor{uavg90}{RGB}{241,243,246}    % a panel's ground
\providecolor{uavg70}{RGB}{214,219,227}    % a gridline
\providecolor{uavg50}{RGB}{179,186,196}    % a reference series
\providecolor{uavg30}{RGB}{116,124,136}    % secondary type
\providecolor{uavg10}{RGB}{45,55,72}       % primary type

% the accents, with the meaning each one carries
\providecolor{uavaccent}{RGB}{220,38,38}    % a bound that is not met
\providecolor{uavwarn}{RGB}{234,120,12}     % a criterion that disagrees
\providecolor{uavok}{RGB}{22,163,74}       % an effect that is resolved


% ---------------------------------------------------------------------
% 2. the axis
%    one spine on the left and one at the bottom, no frame, light dashed
%    gridlines, small dark tick labels and a bold axis label.
% ---------------------------------------------------------------------
\pgfplotsset{
  uavaxis/.style={
      % width/height are the PLOT AREA, so the picture width is a linear
      % function of the axis width and the fit to the column is exact
      scale only axis,
      axis x line=bottom,
      axis y line=left,
      axis line style={uavink, line width=0.6pt},
      tick style={uavink, line width=0.6pt},
      major tick length=2.2pt,
      tick align=outside,
      tick label style={font=\scriptsize, text=uavink, inner sep=1.6pt},
      label style={font=\small\bfseries, text=uavink},
      grid style={uavgrid, dashed, line width=0.5pt},
      legend cell align=left,
      clip=false,
  },
  % the boxed legend strip that runs across the top of a figure
  uavlegend/.style={
      legend style={
          draw=uavgrid,
          fill=uavwhite,
          line width=0.5pt,
          rounded corners=1pt,
          font=\scriptsize,
          text=uavink,
          inner xsep=4pt,
          inner ysep=3pt,
          nodes={inner sep=1.5pt, text=uavink},
          cells={anchor=west},
          column sep=5pt,
          row sep=1pt,
          at={(0.5,1.045)},
          anchor=south,
          /tikz/every even column/.append style={column sep=7pt},
      },
      legend image post style={xscale=0.85},
  },
  % a light grey band: a reference region of the x axis
  uavband/.style={
      draw=none,
      fill=uavband,
      forget plot,
  },
  % a level line: the reference value the figure is read against
  uavlevel/.style={
      draw=uavink,
      line width=0.9pt,
  },
  % the interval bar of a forest plot: thick, round caps, no end marks,
  % a filled marker on the point estimate.  The direction keys cannot be
  % written with "error bars/.cd" inside a style (the path would be
  % changed for every later key of the plot), so they take the full path.
  uavinterval/.style={
      only marks,
      mark=*,
      mark size=2.4pt,
      mark options={line width=0.5pt},
      error bars/error bar style={line width=2.6pt, line cap=round},
      error bars/error mark=none,
  },
  uavintervalx/.style={
      uavinterval,
      error bars/x dir=both,
      error bars/x explicit,
  },
  uavintervaly/.style={
      uavinterval,
      error bars/y dir=both,
      error bars/y explicit,
  },
  % the primary series as a bar: filled, hatched in white, blue edge
  uavhatchbar/.style={
      ybar,
      bar width=0.62,
      fill=uavblue,
      draw=uavblue,
      line width=0.5pt,
      pattern=north east lines,
      pattern color=uavwhite,
  },
  % the reference series as a bar: solid grey, no hatch
  uavrefbar/.style={
      ybar,
      bar width=0.62,
      fill=uavref,
      draw=uavref!70!black,
      line width=0.5pt,
  },
}

\tikzset{
  % the separator between two groups of a categorical axis.  It is used
  % on \draw, which reads keys from the /tikz path, so it is registered
  % with \tikzset rather than in the pgfplots-only namespace above;
  % otherwise pgfkeys reports "/tikz/uavsep unknown" and the paper's
  % gate fails.
  uavsep/.style={
      draw=uavgrid,
      line width=0.5pt,
      dash pattern=on 2pt off 2pt,
  },
}


% ---------------------------------------------------------------------
% 3. the panel label: "(a) By skill", bold, above its panel
% ---------------------------------------------------------------------
\newcommand{\uavpanel}[2]{{\normalfont\bfseries\small (#1)~#2}}


% ---------------------------------------------------------------------
% 4. the table
%    booktabs rules in the theme's colours, and a tinted head row.
%    The colouring is skipped when colortbl is not loaded, so the same
%    macros are safe in a document that only loads plain xcolor.
% ---------------------------------------------------------------------
\newcommand{\uavtoprule}{%
    \ifdefined\arrayrulecolor\arrayrulecolor{uavink}\fi
    \specialrule{1pt}{0pt}{3pt}}
\newcommand{\uavmidrule}{%
    \ifdefined\arrayrulecolor\arrayrulecolor{uavgrid}\fi
    \specialrule{0.8pt}{2pt}{3pt}}
\newcommand{\uavbottomrule}{%
    \ifdefined\arrayrulecolor\arrayrulecolor{uavink}\fi
    \specialrule{1pt}{4pt}{0pt}}
\newcommand{\uavheadrow}{\ifdefined\rowcolor\rowcolor{uavlight}\fi}
\newcommand{\uavzebrarow}{\ifdefined\rowcolor\rowcolor{uavzebra}\fi}
